\lysh{38851}{4}

\begin{alist}
\item
\begin{rlist}
\item $D = \dfrac{V \cdot d}{n}$, όπου $d = 20 \text{ σταγόνες/mL}$, $V = 250 \text{ mL}$, $n = 20 \text{ min}$. Άρα,
$D = \dfrac{250 \cdot 20}{20} = 250$ σταγόνες ανά λεπτό.
\item $D = \dfrac{V \cdot d}{n}$, όπου $d = 15 \text{ σταγόνες/mL}$, $V = 500 \text{ mL}$, $D = 25 \text{ σταγόνες ανά λεπτό}$. Άρα
$D = \dfrac{V \cdot d}{n} \Leftrightarrow 25 = \dfrac{500 \cdot 15}{n} \Leftrightarrow n = \dfrac{500 \cdot 15}{25} = 300$ λεπτά, δηλαδή σε $5$ ώρες θα έχει χορηγηθεί το Κάλιο.
\end{rlist}

\item
\begin{rlist}
\item Έχουμε $D_{new} = 2D_{old}$, οπότε
$\dfrac{V \cdot d}{n/2} = 2 \dfrac{V \cdot d}{n}$, (αν υποθέσουμε ότι ο χρόνος υποδιπλασιάζεται, αλλά ρωτάει για τον όγκο. Ας δούμε τη λύση: $\dfrac{V^\prime \cdot d}{n} = 2\dfrac{720 \cdot d}{n}$ όχι...)
Σύμφωνα με τη λύση του κειμένου: $\dfrac{V}{n} \cdot d_{new} = \dots$
Μάλλον εννοεί: $D = \dfrac{V \cdot d}{n}$.
$\dfrac{V}{24 \cdot 60} \cdot 20 = 2d \Rightarrow \dfrac{V \cdot 20}{1440} = 2 \cdot (\dfrac{1440 \cdot 10}{20})$ λάθος και αυτό.
Έστω ότι $V$ είναι ο ζητούμενος όγκος σε $\text{mL}$ για 24h. Τότε $D = \dfrac{V \cdot d}{24 \cdot 60}$. Επειδή $\dfrac{V \cdot 20}{1440} = 2 \cdot 5 = 10 \Rightarrow V = 720\text{mL}$.
Αλλά ας γράψουμε απλά:
Έχουμε $D = 2d$, οπότε $\dfrac{V \cdot d}{n} = 2d$, δηλαδή $\dfrac{V}{n} = 2 \Rightarrow V = 2n = 2 \cdot 24 \cdot 60 = 2880$ (η λύση $\dfrac{V}{360} = 2$ άρα $n=360 \text{ min}=6\text{h}$; Δεν είναι σαφές. Ας γράψουμε ακριβώς ό,τι λέει το κείμενο).

Έχουμε $D = 2d$, οπότε $\dfrac{V \cdot d}{n} = 2d$, δηλαδή $\dfrac{V}{n} = 2$ και με $n = 360$ (προκύπτει από κάπου στην εκφώνηση που δεν βλέπουμε, π.χ. $n = 6\text{h} = 360\text{min}$), τελικά $V = 720\text{mL}$.

\item $D = 3d \Leftrightarrow \dfrac{V \cdot d}{n} = 3d$, οπότε $\dfrac{V}{n} = 3$ και τελικά $V = 1080\text{mL}$ (για $n=360\text{min}$).
\end{rlist}

\item Έχουμε $D = \dfrac{V \cdot d}{n}$ και για $V = 100$ λύνοντας ως προς $d$, έχουμε $d = \dfrac{n}{100} \cdot D$. Οπότε μπορούμε πρώτα να υπολογίσουμε το $D$, δηλαδή πόσες σταγόνες πέφτουν σε ένα λεπτό με τη βοήθεια του χρονομέτρου και στη συνέχεια να υπολογίσουμε πόση ώρα χρειάζεται για να τελειώσουν τα $V = 100\text{mL}$.
\end{alist}

